%*******************************************************
% Comment lire ce memoire
% Chapitre non numerote, place apres la table des matieres.
%*******************************************************
\pdfbookmark[1]{Comment lire ce mémoire}{CommentLire}
\addcontentsline{toc}{chapter}{\tocEntry{Comment lire ce mémoire}}

\chapter*{Comment lire ce mémoire}
\label{ch:commentlire}
\markboth{Comment lire ce mémoire}{Comment lire ce mémoire}

Ce mémoire s'adresse à des lecteurs venus d'horizons différents. Le traiteur de
signal RADAR ou SONAR connaît déjà les modèles elliptiques et les estimateurs
robustes, mais pas nécessairement le vocabulaire riemannien; le géomètre ou le
praticien de l'apprentissage sur variétés se trouve dans la situation inverse.
La \cref{pt:context}, qui rassemble quatre chapitres d'état des lieux, n'est
donc pas un passage obligé: elle est là pour que chacun puisse combler ce qui
lui manque, et rien de plus. Les contributions elles-mêmes sont regroupées dans
la \cref{pt:contributions}.

Ce chapitre propose quatre choses pour s'y retrouver: une carte des dépendances
entre les chapitres, quelques parcours de lecture selon le profil du lecteur, un
rappel des conventions typographiques et de reproductibilité employées dans tout
le document, et enfin la table des notations.

\section*{Carte des dépendances}
\label{sec:commentlire-carte}

La \cref{fig:commentlire-carte} donne, pour chaque chapitre de contribution ou
de perspective, les sections de l'état des lieux qu'il mobilise effectivement.
Elle se lit de gauche à droite: à gauche les quatre chapitres de mise à niveau
avec le détail de leurs sections, à droite les chapitres qui s'en servent, et
sur chaque flèche les outils empruntés. La couleur désigne la destination, les
traits discontinus les pistes ouvertes du chapitre
\hyperref[ch:perspectives]{\textsc{iii}.4}, qui ne sont
pas des résultats acquis mais des directions de travail.

Une lecture verticale de la carte est instructive: le chapitre
\hyperref[ch:spdnet]{\textsc{iii}.3} est le seul à mobiliser les quatre
chapitres de l'état des lieux, et il le fait pour des raisons de nature
différente --- le régime dimensionnel pour former ses matrices, la loi de
Wishart et les divergences pour choisir ses moyennes, la variété de Stiefel pour
contraindre ses poids. Une lecture horizontale l'est tout autant: la
\cref{sec:context-riemann-optimization} sur l'optimisation riemannienne est le
seul outil que les trois chapitres de contribution empruntent tous les trois.

Le chapitre \hyperref[ch:learning]{\textsc{iii}.2} fait exception à la logique
d'ensemble, et il faut le signaler ici. Il a besoin de deux notions --- les
modèles graphiques gaussiens et la théorie des matrices aléatoires --- qui ne
figurent nulle part dans la \cref{pt:context}, et qui sont donc introduites dans
le chapitre lui-même, au strict nécessaire. Ce n'est pas un oubli: la partie
d'état des lieux n'a vocation à contenir que ce qui sert à plusieurs chapitres,
et ces deux notions-là ne servent qu'une fois.

Deux remarques sur l'usage de cette carte. D'abord, une flèche signale une
dépendance réelle, pas un prérequis absolu: un lecteur pressé peut toujours
accepter un outil sur parole et revenir plus tard sur sa construction. Ensuite,
l'absence de flèche est aussi une information: la \cref{sec:context-geometry-icrb}
sur la borne de Cramér-Rao intrinsèque, par exemple, n'est utilisée par aucune
contribution de ce mémoire; elle est là parce qu'elle éclaire ce que mesure la
métrique de Fisher, et peut être sautée sans dommage.

\begin{figure}[p]
  \centering
  % la carte deborde legerement dans les marges pour rester lisible
  \makebox[\textwidth][c]{\scalebox{1.06}{% Carte des dependances du memoire: chapitres de l'etat des lieux (colonne de
% gauche) vers les chapitres de contributions et de perspectives (colonne de
% droite). Chaque arete porte les sections effectivement mobilisees; la couleur
% designe la destination. Utilise dans FrontBackmatter/CommentLire.tex.
%
% Revision du 2026-08-27: insertion du \cref{ch:learning} entre le chapitre de
% detection et celui sur SPDnet. Il mobilise trois des quatre chapitres de
% l'etat des lieux --- pas le chapitre 4, ce qui laisse au \cref{ch:spdnet} son
% statut de seul chapitre a les mobiliser tous les quatre. Les deux notions
% nouvelles dont il a besoin (modeles graphiques, matrices aleatoires) sont
% introduites dans le chapitre lui-meme et n'apparaissent donc pas ici.
%
% Revision du 2026-08-24: le \cref{ch:spdnet} etant ecrit, ses dependances
% reelles sont desormais connues et elles sont bien plus larges que la seule
% arete depuis le chapitre 3 qui figurait ici. Il mobilise en fait les quatre
% chapitres de l'etat des lieux --- c'est meme le seul chapitre du memoire dans
% ce cas. Trois aretes ont donc ete ajoutees et une entree (§2.1.1, Wishart) a
% ete ouverte dans le bloc du chapitre 2.
\definecolor{lctx}{RGB}{90,90,90}
\definecolor{lcinq}{RGB}{40,70,140}
\definecolor{lsix}{RGB}{110,80,140}
\definecolor{lsept}{RGB}{50,110,120}
\definecolor{lhuit}{RGB}{125,125,125}
\begin{tikzpicture}[
    >={Stealth[length=4pt,width=3pt]},
    line join=round,
    font=\footnotesize,
    entry/.style={anchor=west, inner sep=1pt, font=\scriptsize, text depth=0.5ex},
    ctitle/.style={anchor=west, inner sep=1pt, font=\scriptsize\bfseries, text=lctx},
    frame/.style={draw=lctx!45, rounded corners=2pt, inner xsep=5pt, inner ysep=6pt},
    dest/.style={draw, rounded corners=2pt, align=center, font=\scriptsize,
                 text width=2.5cm, inner sep=5pt},
    elab/.style={font=\tiny, fill=white, inner sep=1.5pt, align=center},
    edge/.style={thin, ->},
  ]

  %--------------------------------------------------------------- etat des lieux
  % Les numeros de chapitre repartent a 1 dans chaque partie
  % (\counterwithin*{chapter}{part}): on les prefixe donc par la partie.
  \node[font=\scriptsize\bfseries, text=lctx, anchor=west] at (0,1.35)
    {Partie \textsc{ii} --- État des lieux};
  \node[ctitle] (t1) at (0,0.62) {ch.\,1 --- Covariance};
  \node[entry]  (a1) at (0, 0.00) {\S1.3 \ estimation, SCM};
  \node[entry]  (a2) at (0,-0.55) {\S1.4 \ régularisation};
  \node[entry]  (a3) at (0,-1.10) {\S1.5 \ cas complexe};

  \node[ctitle] (t2) at (0,-2.05) {ch.\,2 --- Distributions};
  \node[entry]  (b0) at (0,-2.67) {\S2.1.1 gaussienne, Wishart};
  \node[entry]  (b1) at (0,-3.22) {\S2.1.2 elliptiques};
  \node[entry]  (b2) at (0,-3.77) {\S2.1.3 M-estimation};
  \node[entry]  (b3) at (0,-4.32) {\S2.2 \ cas complexe};

  \node[ctitle] (t3) at (0,-5.27) {ch.\,3 --- Riemann};
  \node[entry]  (c1) at (0,-5.89) {\S3.1 \ variétés, transport};
  \node[entry]  (c2) at (0,-6.44) {\S3.2 \ cône SPD (AI, LE)};
  \node[entry]  (c3) at (0,-6.99) {\S3.3 \ moyennes de Fréchet};
  \node[entry]  (c4) at (0,-7.54) {\S3.4 \ matrices orthog.};
  \node[entry]  (c5) at (0,-8.09) {\S3.5 \ optimisation};

  \node[ctitle] (t4) at (0,-9.04) {ch.\,4 --- Géométrie de l'information};
  \node[entry]  (d1) at (0,-9.66) {\S4.1 \ métrique de Fisher};
  \node[entry]  (d2) at (0,-10.21) {\S4.2 \ modèle gaussien};
  \node[entry]  (d3) at (0,-10.76) {\S4.5 \ BCR intrinsèque};
  \node[entry]  (d4) at (0,-11.31) {\S4.6 \ divergences};

  \coordinate (r) at (4.62,0); % garantit une largeur commune aux quatre cadres

  \begin{scope}[on background layer]
    \node[frame, fit=(t1)(a1)(a3)(r)]        (F1) {};
    \node[frame, fit=(t2)(b0)(b3)(r|-b0)]    (F2) {};
    \node[frame, fit=(t3)(c1)(c5)(r|-c1)]    (F3) {};
    \node[frame, fit=(t4)(d1)(d4)(r|-d1)]    (F4) {};
  \end{scope}

  %--------------------------------------------------------------- destinations
  \node[dest, draw=lcinq, text=lcinq, fill=lcinq!5] (ch5) at (10.9,-2.10)
    {\textbf{\textsc{iii}.1}\\Test d'hypothèses robustes\\\textit{sonar, RSO}};
  \node[dest, draw=lsix, text=lsix, fill=lsix!5] (ch6) at (10.9,-4.35)
    {\textbf{\textsc{iii}.2}\\Apprendre sur\\des covariances};
  \node[dest, draw=lsept, text=lsept, fill=lsept!5] (ch7) at (10.9,-6.90)
    {\textbf{\textsc{iii}.3}\\Apprentissage profond\\sur la variété SPD};
  \node[dest, draw=lhuit, dashed, text=lhuit] (ch8a) at (10.9,-9.55)
    {\textbf{\textsc{iii}.4}, \S2\\Transformers SPD};
  \node[dest, draw=lhuit, dashed, text=lhuit] (ch8b) at (10.9,-11.05)
    {\textbf{\textsc{iii}.4}, \S3\\Phase géométrique};

  \node[font=\scriptsize\bfseries, text=lhuit, anchor=south west] at (9.45,-8.75)
    {Perspectives};
  \node[font=\scriptsize\bfseries, text=lcinq!80, anchor=south] at (10.9,-1.05)
    {Contributions};

  %--------------------------------------------------------------- aretes
  % ---- vers le chapitre III.1
  \begin{scope}[edge, lcinq]
    \draw (F1.east) to[out=0,in=165]
      node[elab,pos=0.52,text=lcinq]{\S1.3 SCM \textperiodcentered\ \S1.4 régularisation\\\S1.5 circularité}
      (ch5.west);
    \draw ($(F2.east)+(0,0.30)$) to[out=0,in=195]
      node[elab,pos=0.62,text=lcinq]{\S2.1.2 SIRV\\\S2.1.3 Tyler, $M$-estimation}
      (ch5.west);
    \draw ($(F3.east)+(0,0.95)$) to[out=25,in=215]
      node[elab,pos=0.32,yshift=-13pt,text=lcinq]{\S3.2 distance affine invariante\\\S3.3 moyennes \textperiodcentered\ \S3.5 descente}
      (ch5.south);
    % contournement par la droite: ne pas traverser le bloc III.2
    \draw[rounded corners=4pt] ($(F4.east)+(0,0.55)$)
      -- (9.05,-8.49) -- (9.05,-2.90)
      node[elab,pos=0.86,text=lcinq]{\S4.1 test de Rao\\\S4.6 divergences}
      -- ($(ch5.west)+(0,-0.45)$);
  \end{scope}

  % ---- vers le chapitre III.2 (apprentissage sur des covariances)
  % Trois aretes seulement: ce chapitre n'emprunte rien au chapitre 4, et les
  % deux notions dont il a besoin en propre (modeles graphiques, matrices
  % aleatoires) sont introduites dans le chapitre lui-meme.
  \begin{scope}[edge, lsix]
    % ch.1: le regime dimensionnel, dont la theorie des matrices aleatoires est
    % la description fine.
    \draw ($(F1.east)+(0,-0.75)$) to[out=-25,in=170]
      node[elab,pos=0.25,text=lsix]{\S1.4 régime\\dimensionnel}
      ($(ch6.west)+(0,0.30)$);
    % ch.2: la vraisemblance elliptique penalisee de l'apprentissage de graphes.
    \draw ($(F2.east)+(0,-0.95)$) to[out=-12,in=178]
      node[elab,pos=0.20,yshift=-3pt,text=lsix]{\S2.1.2 elliptiques\\\S2.1.3 $M$-estimation}
      (ch6.west);
    % ch.3: cone, moyennes et descente riemannienne, communs aux deux travaux.
    \draw ($(F3.east)+(0,1.30)$) to[out=18,in=200]
      node[elab,pos=0.74,yshift=11pt,text=lsix]{\S3.2 cône \textperiodcentered\ \S3.3 moyennes\\\S3.5 optimisation}
      ($(ch6.west)+(0,-0.32)$);
  \end{scope}

  % ---- vers le chapitre III.3 (apprentissage profond sur la variete SPD)
  % Quatre aretes: c'est le seul chapitre du memoire a mobiliser les quatre
  % chapitres de l'etat des lieux.
  \begin{scope}[edge, lsept]
    % ch.1: la covariance apprise est une SCM, et son conditionnement decide
    % du seuil de ReEig. Contournement par le haut et par la droite.
    \draw[rounded corners=4pt] ($(F1.east)+(0,0.20)$)
      -- (12.85,0.20)
      node[elab,pos=0.55,text=lsept]{\S1.3 covariance empirique\\\S1.4 conditionnement}
      -- (12.85,-6.90) -- (ch7.east);
    % ch.2: Wishart et Wishart inverse donnent la lecture statistique des
    % moyennes arithmetique et harmonique.
    \draw ($(F2.east)+(0,-0.45)$) to[out=-20,in=150]
      node[elab,pos=0.82,text=lsept]{\S2.1.1 Wishart} ($(ch7.west)+(0,0.42)$);
    % ch.3: le gros de la dependance.
    \draw ($(F3.east)+(0,0.05)$) to[out=0,in=180]
      node[elab,pos=0.48,text=lsept]{\S3.2 cône (AI, LE) \textperiodcentered\ \S3.3 moyennes\\\S3.4 Stiefel \textperiodcentered\ \S3.5 rétractions}
      (ch7.west);
    % ch.4: c'est lui qui fait de la liste des moyennes une famille indexee par
    % un modele statistique.
    \draw ($(F4.east)+(0,-0.35)$) to[out=10,in=250]
      node[elab,pos=0.30,text=lsept]{\S4.2 distance de Rao gaussienne\\\S4.6 divergences}
      ($(ch7.south)+(0.35,0)$);
  \end{scope}

  % ---- vers les perspectives
  \begin{scope}[edge, lhuit, dashed]
    \draw ($(F3.east)+(0,-0.85)$) to[out=-15,in=180]
      node[elab,pos=0.62,text=lhuit]{\S3.3 moyenne de\\Karcher pondérée} (ch8a.west);
    \draw ($(F3.east)+(0,-1.35)$) to[out=-30,in=180]
      node[elab,pos=0.82,text=lhuit]{\S3.1 transport\\parallèle} (ch8b.west);
    \draw[rounded corners=4pt] ($(F1.east)+(0,-0.40)$)
      -- (5.25,-0.95) -- (5.25,-12.05)
      -- (10.9,-12.05)
      node[elab,pos=0.88,text=lhuit]{\S1.5 pseudo-covariance} (ch8b.south);
    \draw ($(ch7.south east)+(-0.90,0)$) -- ($(ch8a.north east)+(-0.90,0)$);
  \end{scope}

  %--------------------------------------------------------------- legende
  \node[anchor=north west, font=\tiny, text=lctx, align=left] at (0,-12.40)
    {\textcolor{lcinq}{\rule[0.4ex]{7pt}{0.5pt}}~outil mobilisé en \textsc{iii}.1
     \quad
     \textcolor{lsix}{\rule[0.4ex]{7pt}{0.5pt}}~en \textsc{iii}.2
     \quad
     \textcolor{lsept}{\rule[0.4ex]{7pt}{0.5pt}}~en \textsc{iii}.3
     \quad
     \textcolor{lhuit}{\rule[0.4ex]{3pt}{0.5pt}\hspace{1.5pt}\rule[0.4ex]{3pt}{0.5pt}}~piste ouverte en \textsc{iii}.4
     \\
     (les numéros de chapitre repartent à 1 dans chaque partie)};
\end{tikzpicture}
}}
  \caption[Carte des dépendances du mémoire]{Carte des dépendances. Les
    chapitres de l'état des lieux (\cref{pt:context}) sont à gauche avec le
    détail de leurs sections; les chapitres de contribution
    (\cref{pt:contributions}) et de perspective sont à droite. Chaque flèche
    porte les sections effectivement mobilisées, et sa couleur désigne sa
    destination. Les traits discontinus mènent aux pistes du chapitre
    \textsc{iii}.4.}
  \label{fig:commentlire-carte}
\end{figure}

\section*{Parcours de lecture}
\label{sec:commentlire-parcours}

Cinq entrées possibles, selon ce que l'on connaît déjà et ce que l'on cherche.

\begin{figure}[htb]
  \centering
  \makebox[\textwidth][c]{% Les quatre parcours de lecture, en regard de la section « Parcours de
% lecture » de FrontBackmatter/CommentLire.tex. Meme palette que carte.tex:
% chaque couloir est une suite de jalons, du point d'entree au chapitre vise.
\definecolor{lctx}{RGB}{90,90,90}
\definecolor{lcinq}{RGB}{40,70,140}
\definecolor{lsix}{RGB}{50,110,120}
\definecolor{lsept}{RGB}{125,125,125}
\begin{tikzpicture}[
    >={Stealth[length=4pt,width=3pt]},
    line join=round,
    font=\footnotesize,
    lane/.style={anchor=east, align=right, text width=2.3cm,
                 font=\scriptsize\bfseries, inner sep=0pt},
    jalon/.style={draw, rounded corners=2pt, align=center, font=\scriptsize,
                 text width=1.75cm, inner sep=4pt, minimum height=0.95cm},
    goal/.style={jalon, very thick},
    hop/.style={thin, ->, shorten >=1pt, shorten <=1pt},
    note/.style={font=\tiny, text=lctx, align=center},
  ]

  % abscisses des quatre colonnes de jalons
  \def\ca{3.45} \def\cb{6.05} \def\cc{8.65} \def\cd{11.25}

  %--------------------------------------------------- 1. radar / sonar
  \node[lane, text=lcinq] at (2.30,0) {Traitement du signal RADAR ou SONAR};
  \node[jalon, draw=lcinq!55, text=lcinq] (a1) at (\ca,0) {bilans des ch.~\textsc{ii}.1 et \textsc{ii}.2};
  \node[jalon, draw=lcinq!55, text=lcinq] (a2) at (\cb,0) {\S3.2\\cône SPD};
  \node[jalon, draw=lcinq!55, text=lcinq] (a3) at (\cc,0) {\S3.3\\moyennes};
  \node[goal, draw=lcinq, text=lcinq, fill=lcinq!7] (a4) at (\cd,0) {ch.~\textsc{iii}.1\\détection};
  \draw[hop, lcinq] (a1) -- (a2);
  \draw[hop, lcinq] (a2) -- (a3);
  \draw[hop, lcinq] (a3) -- (a4);
  \node[note, anchor=north] at ($(a3)!0.5!(a4)+(0,-0.50)$)
    {\S4.1 en appoint\\(tests de Rao)};

  %--------------------------------------------------- 2. geometres
  \node[lane, text=lsix] at (2.30,-2.05) {Géométrie riemannienne, apprentissage sur variétés};
  \node[jalon, draw=lsix!55, text=lsix] (b1) at (\ca,-2.05) {\S2.1.2\\elliptiques};
  \node[jalon, draw=lsix!55, text=lsix] (b2) at (\cb,-2.05) {\S2.1.3\\M-estimation};
  \node[jalon, draw=lsix!55, text=lsix] (b3) at (\cc,-2.05) {ch.~\textsc{ii}.4\\Fisher};
  \node[goal, draw=lsix, text=lsix, fill=lsix!7] (b4) at (\cd,-2.05) {ch.~\textsc{iii}.2 et \textsc{iii}.3};
  \draw[hop, lsix] (b1) -- (b2);
  \draw[hop, lsix] (b2) -- (b3);
  \draw[hop, lsix] (b3) -- (b4);

  %--------------------------------------------------- 3. lecture rapide
  \node[lane, text=lsept] at (2.30,-4.10) {Lecture rapide};
  \node[jalon, draw=lsept!60, text=lsept] (c1) at (\ca,-4.10) {curriculum vitae};
  \node[jalon, draw=lsept!60, text=lsept] (c2) at (\cb,-4.10) {bilans des ch.~\textsc{ii}.1 à \textsc{ii}.4};
  \node[goal, draw=lsept, text=lsept, fill=lsept!7] (c3) at (\cc,-4.10)
    {intro.\ et bilans de la partie \textsc{iii}};
  \draw[hop, lsept] (c1) -- (c2);
  \draw[hop, lsept] (c2) -- (c3);
  \node[note, anchor=west] at ($(c3.east)+(0.35,0)$)
    {les blocs \emph{Contributions}\\donnent les publications};

  %--------------------------------------------------- 4. lecture lineaire
  \node[lane, text=lctx] at (2.30,-5.85) {Lecture linéaire};
  \node[jalon, draw=lctx!55, text=lctx, text width=8.65cm, minimum height=0.7cm,
        anchor=west] (d1) at (\ca-0.95,-5.85)
    {l'intégralité, dans l'ordre du mémoire: chaque chapitre de l'état des lieux
     naît d'une insuffisance du précédent};

  %--------------------------------------------------- reperes de colonne
  \node[note, anchor=south] at (\ca,0.75) {point d'entrée};
  \node[note, anchor=south] at (\cd,0.75) {chapitre visé};
  \draw[lctx!30, dashed, thin] (\ca,0.65) -- (\ca,-5.35);
  \draw[lctx!30, dashed, thin] (\cd,0.65) -- (\cd,-2.75);
\end{tikzpicture}
}
  \caption[Les parcours de lecture]{Les parcours de lecture.}
  \label{fig:commentlire-parcours}
\end{figure}

\paragraph{Traitement du signal RADAR ou SONAR.} Le contenu des chapitres
\hyperref[ch:context-definition]{\textsc{ii}.1} et
\hyperref[ch:context-distributions]{\textsc{ii}.2} est probablement acquis:
on pourra n'en lire que les bilans, puis prendre la
\cref{sec:context-riemann-spd} sur le cône SPD et la
\cref{sec:context-riemann-moyennes} sur les moyennes, qui sont les deux seules
véritables nouveautés de vocabulaire nécessaires, et enchaîner directement sur
le chapitre \hyperref[ch:robust-detection]{\textsc{iii}.1}. La
\cref{sec:context-geometry-fisher} rendra
service au moment des tests de Rao.

\paragraph{Géométrie riemannienne et apprentissage sur variétés.} La situation
symétrique: le chapitre \hyperref[ch:context-riemann]{\textsc{ii}.3} sera du
déjà-vu, mais la
\cref{sec:context-distributions-reel} y est indispensable, en particulier ses
sous-sections sur les distributions elliptiques et la M-estimation, qui portent
toute la modélisation statistique du chapitre
\hyperref[ch:robust-detection]{\textsc{iii}.1}. Pour le chapitre
\hyperref[ch:spdnet]{\textsc{iii}.3}, ce sont plutôt la loi de Wishart
(\cref{def:context-distributions-wishart}) et les divergences de la
\cref{sec:context-geometry-divergences} qu'il faut avoir en tête: ce sont elles
qui font des différentes moyennes du cône une famille indexée par un modèle
statistique, et non un catalogue de recettes. Le chapitre
\hyperref[ch:context-geometry]{\textsc{ii}.4} fait le pont entre les deux mondes et
constitue, de ce point de vue, le chapitre le plus rentable du mémoire.

\paragraph{Apprentissage statistique.} Un troisième profil, apparu avec le
chapitre \hyperref[ch:learning]{\textsc{iii}.2}: celui qui vient de
l'apprentissage et cherche ce que la modélisation statistique lui apporte. Le
plus court chemin est la \cref{sec:context-covariance-regularized-estimation}
sur le régime dimensionnel, puis la \cref{sec:context-riemann-spd} sur le cône
et la \cref{sec:context-riemann-moyennes} sur les moyennes, et enfin la
\cref{sec:context-riemann-optimization} sur l'optimisation riemannienne --- qui
est l'outil commun aux trois chapitres de contribution. Les modèles graphiques
et la théorie des matrices aléatoires dont ce chapitre a besoin sont introduits
sur place et ne demandent aucun détour.

\paragraph{Lecture rapide.} Pour prendre la mesure de l'ensemble sans entrer
dans le détail: le \hyperref[ch:cv]{curriculum vitae}, puis la section de bilan qui clôt
chacun des quatre chapitres de l'état des lieux, puis l'introduction et le bilan
de chaque chapitre de contribution. Les blocs \textit{Contributions}, signalés
en marge à l'ouverture des chapitres de la \cref{pt:contributions}, résument en
un paragraphe les travaux dont chaque chapitre rend compte et les publications
correspondantes.

\paragraph{Lecture linéaire.} Enfin, pour qui découvre le sujet, l'ordre du
mémoire est un ordre de construction: chaque chapitre de l'état des lieux naît
d'une insuffisance du précédent, et la \cref{pt:contributions} ne fait
qu'assembler ce qui a été mis en place. C'est le parcours le plus long, et le
seul où les renvois croisés se lisent sans effort.

\section*{Conventions}
\label{sec:commentlire-conventions}

\paragraph{Numérotation.} Les chapitres repartent à~1 dans chaque partie, ainsi
que les figures, les équations et les blocs numérotés qui en dérivent. Deux
chapitres de parties différentes portent donc le même numéro, et lorsque le
contexte ne suffit pas à lever l'ambiguïté --- c'est le cas dans la carte
ci-dessus --- le numéro de chapitre est préfixé par celui de la partie:
\textsc{ii}.3 désigne le troisième chapitre de l'état des lieux,
\textsc{iii}.1 le premier chapitre de contributions.

\paragraph{Notations.} Les vecteurs et les matrices sont en gras
($\x$, $\MCov$), les scalaires en italique, les ensembles et variétés en
calligraphié ($\Manifold$). Sauf mention contraire, tout est écrit dans le
\textbf{cas réel}: c'est le cas où les dimensions et les figures se lisent le
plus facilement. Chaque chapitre de l'état des lieux se termine par une section
consacrée au cas complexe, qui établit le dictionnaire vers les données
effectivement traitées en RADAR, en SONAR et en RSO; les chapitres de
contribution, eux, se placent d'emblée dans le cas complexe. La table des
notations qui clôt ce chapitre rassemble l'ensemble des symboles employés.

\paragraph{Marges.} Les marges de ce document sont larges et servent
activement. Les références bibliographiques y apparaissent en regard de l'appel
de citation, de sorte que l'on puisse identifier un travail cité sans aller à la
bibliographie. Les blocs numérotés --- définitions, propositions, remarques ---
portent leur intitulé en marge plutôt qu'en tête, ce qui permet de les repérer
en feuilletant. Les notes de bas de page restent en bas de page.

\paragraph{Figures reproductibles.} Toutes les figures de simulation de ce
mémoire sont produites par des expériences numériques versionnées, et chacune
porte en légende, en gris et en caractères à chasse fixe, la provenance exacte
de ce qui est tracé. Ainsi, la \cref{fig:context-riemann-moyennes} affiche:
\begin{center}
  \provMoyennes
\end{center}
soit, dans l'ordre: le \texttt{sha} du commit du dépôt de simulations qui a
produit la figure, la date d'exécution, la graine du générateur aléatoire, puis
la valeur de \emph{tous} les paramètres de l'expérience --- ici quinze matrices
tirées avec une dispersion de $0{,}9$ et un conditionnement de $4$, une moyenne
calculée à $10^{-10}$ près en au plus cent itérations. Ces lignes sont
volontairement exhaustives: elles suffisent, à elles seules, à rejouer la
figure.

Le QR code placé en marge de la même figure pointe vers la page de l'expérience
correspondante, où l'on trouve le code source, les paramètres et les données
produites\footnote{À l'adresse
\url{https://ammarmian.fr/simulations_hdr/}.}. La chaîne complète est la
suivante: chaque expérience est décrite par un fichier de configuration et
exécutée sous \texttt{qanat}, qui enregistre le commit, les paramètres et les
sorties; l'exécution exporte la figure en PGFPlots ainsi qu'un fichier de
provenance; celui-ci est enfin traduit en commandes \LaTeX{} que la légende
appelle. La ligne grise n'est donc pas recopiée à la main: elle est engendrée en
même temps que la figure, et ne peut pas se désynchroniser d'elle.

Les figures explicatives, elles --- schémas de géométrie, illustrations de
définitions --- sont dessinées directement en Ti\emph{k}Z et ne portent pas de
ligne de provenance: il n'y a rien à reproduire.

\section*{Table des notations}
\label{sec:commentlire-notations}

Les symboles sont rassemblés ici par domaine, dans l'ordre où les chapitres les
introduisent. Sauf mention contraire, ils sont donnés dans le cas réel; le
passage au cas complexe remplace la transposition $(\cdot)^\T$ par la
transposition conjuguée $(\cdot)^\herm$, les matrices symétriques par les
matrices hermitiennes, et l'accent tilde marque les grandeurs correspondantes
lorsqu'elles diffèrent.

\newcommand{\notatable}[1]{%
  \begingroup
  \small
  \setlength{\LTpre}{0.2\baselineskip}\setlength{\LTpost}{1.0\baselineskip}
  \renewcommand{\arraystretch}{1.15}
  \begin{longtable}{@{}p{0.24\textwidth}p{0.72\textwidth}@{}}
  #1
  \end{longtable}
  \endgroup}

\subsection*{Généralités et algèbre linéaire}
\notatable{
  $\R$, $\C$, $\N$ & corps des réels, des complexes, entiers naturels \\
  $\Nvar$ & dimension des vecteurs observés (nombre de variables) \\
  $\Nobs$ & nombre d'observations disponibles \\
  $\Id{\Nvar}$ & matrice identité de taille $\Nvar$ \\
  $(\cdot)^\T$, $(\cdot)^\herm$, $\conj{(\cdot)}$ & transposition, transposition conjuguée, conjugaison \\
  $\inv{(\cdot)}$, $\Det(\cdot)$, $\trace{\cdot}$ & inverse, déterminant, trace \\
  $\Realp(\cdot)$, $\Imagp(\cdot)$ & parties réelle et imaginaire \\
  $\Diag(\cdot)$, $\Cond(\cdot)$ & matrice diagonale, conditionnement \\
  $\norm{\cdot}$, $\norm{\cdot}_\frob$, $\abs{\cdot}$ & norme, norme de Frobenius, module \\
  $\expm(\cdot)$, $\logm(\cdot)$ & exponentielle et logarithme de matrice \\
  $\symop(\cdot)$, $\hermop(\cdot)$ & parties symétrique et hermitienne \\
  $\vecop(\cdot)$, $\Kmat$ & vectorisation, matrice de commutation \\
  $\otimes$, $\had$ & produits de Kronecker et de Hadamard \\
  $\eig$, $\eigset$ & valeur propre, matrice des valeurs propres \\
  $\Sym{\Nvar}$, $\Herm{\Nvar}$ & matrices symétriques réelles, hermitiennes \\
  $\SPD{\Nvar}$, $\HPD{\Nvar}$ & matrices définies positives, réelles et hermitiennes \\
  $\SPDdet{\Nvar}$, $\HPDdet{\Nvar}$ & les mêmes, de déterminant unité \\
  $\Toep{\Nvar}$ & matrices de Toeplitz de taille $\Nvar$ \\
  $\complexform{\cdot}$ & forme réelle d'une matrice complexe \\
  $\eqdef$, $\bigO(\cdot)$ & égalité par définition, ordre de grandeur \\
}

\subsection*{Statistiques et probabilités}
\notatable{
  $\Universe$, $\event$ & univers, événement élémentaire \\
  $\Variable$, $\x$ & vecteur aléatoire, réalisation observée \\
  $\X$, $\Xset$ & données rassemblées en matrice, en ensemble \\
  $\E\{\cdot\}$, $\Prob$ & espérance, probabilité \\
  $\Vmoy$, $\MCov$ & moyenne et matrice de covariance \\
  $\MPCov$ & matrice de pseudo-covariance (cas complexe) \\
  $\Smat$ & matrice des moments d'ordre deux $\X\X^\T$ \\
  $\ddp{\param}$, $\likelihood$ & densité de probabilité, log-vraisemblance \\
  $\param$ & paramètre générique d'un modèle \\
  $\score$, $\Fisher$ & vecteur score, matrice d'information de Fisher \\
  $\gaussian$, $\cgaussian$ & loi gaussienne réelle, complexe circulaire \\
  $\gcgaussian$ & loi gaussienne complexe non circulaire \\
  $\elliptical$, $\celliptical$ & lois elliptiques, réelles et complexes \\
  $\compoundgaussian$, $\ccompoundgaussian$ & lois gaussiennes composées \\
  $\studentdist$, $\kdist$, $\gengaussian$ & lois de Student, $K$ et gaussienne généralisée \\
  $\gammadist$, $\chisquared$, $\uniform$ & lois gamma, du khi-deux, uniforme \\
  $\wishart$, $\cwishart$ & lois de Wishart, réelle et complexe \\
  $\angular$, $\cangular$ & lois angulaires elliptiques \\
}

\subsection*{Distributions elliptiques et estimation robuste}
\notatable{
  $\densgen$ & générateur de densité \\
  $\logderiv$ & dérivée logarithmique $\densgen'/\densgen$ \\
  $\densgenmoment{k}$, $\densgenconst{\Nvar}$ & moments et constante de normalisation du générateur \\
  $\modvar$, $\quadvar$ & variable modulaire, forme quadratique aléatoire \\
  $\sphere{\Nvar-1}$, $\spherearea{\Nvar}$ & sphère unité et son aire \\
  $\texture$, $\dof$ & texture, paramètre de forme (degrés de liberté) \\
  $\MScatter$, $\MShape$ & matrice de dispersion, matrice de forme \\
  $\Mestim$, $\Tyler$ & $M$-estimateur de dispersion, estimateur de Tyler \\
  $\weight$, $\scorefun$ & fonction de poids, fonction $\scorefun(t)=t\,\weight(t)$ \\
  $\coefledoitwolf$, $\targetledoitwolf$ & coefficient et cible d'un estimateur régularisé \\
  $\iiter$, $\Niter$ & indice et nombre d'itérations d'un algorithme \\
}

\subsection*{Géométrie riemannienne}
\notatable{
  $\Manifold$, $\ambient$ & variété, espace ambiant euclidien \\
  $\Tangent{\param}{\Manifold}$ & espace tangent à $\Manifold$ en $\param$ \\
  $\tvec$, $\tvecbis$ & vecteurs tangents \\
  $\metric{\tvec}{\tvecbis}{\param}$ & métrique riemannienne au point $\param$ \\
  $\normriem{\tvec}{\param}$ & norme induite par la métrique \\
  $\projman_\param$ & projection orthogonale sur l'espace tangent \\
  $\connection$, $\transport$ & connexion de Levi-Civita, transport parallèle \\
  $\geod$, $\geodist$ & géodésique, distance géodésique \\
  $\geodmean$ & moyenne géodésique de deux points \\
  $\Fmean$ & moyenne de Fréchet (de Karcher) \\
  $\Retr$ & rétraction \\
  $\Rgrad$, $\Rhess$, $\dirderiv$ & gradient et hessien riemanniens, dérivée directionnelle \\
  $\Stiefel{\Nvar}{\Nrank}$, $\Grassmann{\Nvar}{\Nrank}$ & variétés de Stiefel et de Grassmann \\
  $\Orthgroup{\Nvar}$, $\Unitgroup{\Nvar}$ & groupes orthogonal et unitaire \\
  $\Rpos{\Nvar}$ & cône des vecteurs à entrées strictement positives \\
  $\loss$, $\risk$, $(\cdot)^\opt$ & fonction de coût, risque, point optimal \\
  $\KL$ & divergence de Kullback-Leibler \\
  $\metriccoefa$, $\metriccoefb$ & coefficients de la métrique de Fisher elliptique \\
}

\subsection*{Détection et modèles de signal}
\notatable{
  $\hypothesis_0$, $\hypothesis_1$ & hypothèses nulle et alternative \\
  $\DetectorStat \detect$ & règle de décision au seuil $\DetectorTreshold$ \\
  $\statGLRT$, $\statRao$ & statistiques du \detname{glrt} et du test de Rao \\
  $\statWald$, $\statANMF$ & statistiques du test de Wald et de l'\detname{anmf} \\
  $\pfa$, $\pdet$, $\snr$ & probabilités de fausse alarme, de détection, rapport signal à bruit \\
  $\paramtest$, $\paramnuis$ & paramètres testés et paramètres de nuisance \\
  $\z$, $\speckle$ & bruit observé, composante speckle gaussienne \\
  $\amplitude$, $\Vamplitude$ & amplitude complexe de la cible, vecteur des amplitudes \\
  $\steering$, $\steerfull$ & vecteurs directionnels, par antenne et complet \\
  $\Steermat$, $\Steerspace$, $\Nrank$ & matrice directionnelle, base du sous-espace signal, son rang \\
  $\projspace{\Steerspace}$ & projecteur sur le sous-espace signal \\
  $\Nsensor$, $\Narray$, $\Nping$ & nombres de capteurs, d'antennes et de pings \\
  $\wavelength$, $\doa$ & longueur d'onde centrale, direction d'arrivée \\
  $\powerscale$, $\escale$, $\Scalemat$ & facteur d'échelle, échelle effective, matrice des échelles \\
  $\Kronleft \otimes \Kronright$ & structure de Kronecker de la dispersion \\
  $\dimleft$, $\dimright$ & dimensions des deux opérandes, $\Nvar=\dimleft\dimright$ \\
  $\obsmat$ & remise en matrice d'une observation vectorisée \\
  $\Ntime$, $\itime$ & nombre d'instants et indice temporel \\
  $\HH$, $\HV$, $\VH$, $\VV$ & canaux de polarisation en imagerie RSO \\
}

\subsection*{Apprentissage sur les covariances}
\notatable{
  $\Nmat$, $\imat$ & nombre de covariances de la collection, indice associé \\
  $\SCM$ & covariance empirique (scm) \\
  $\cratio$ & rapport de concentration $\Nvar/\Nobs$ \\
  $\MPdens$, $\eig_\pm$ & densité de Mar\v{c}enko-Pastur et bornes de son support \\
  $\eigvec$, $\Eigmat$ & vecteur et matrice diagonale des valeurs propres \\
  $\distrmt$, $\Fmeanrmt$ & distance et moyenne de Fréchet corrigées par la théorie des matrices aléatoires \\
  $\Diagpos{\Nvar}$ & matrices diagonales définies positives \\
  $\Nclass$, $\iclass$, $\classlabel$ & nombre de classes, indice de classe, étiquette \\
  $\Fmeanclass{\iclass}$ & moyenne de Fréchet de la classe $\iclass$ \\
  $\graph=(\vertexset,\edgeset)$ & graphe appris, ses sommets et ses arêtes \\
  $\Precision=\inv{\MCov}$ & matrice de précision \\
  $\penfun$, $\penaltyelem$, $\regparam$ & pénalité, pénalité élément par élément, paramètre de régularisation \\
  $\Lowrank$, $\Diagpart$, $\Factormat$ & c\oe ur de rang faible, partie diagonale, matrice des facteurs \\
  $\PSDrank{\Nvar}{\Nrank}$ & matrices semi-définies positives de rang $\Nrank$ \\
  $\Factorman{\Nvar}{\Nrank}$ & matrices « rang faible plus diagonale » \\
}
